problem:  
Let \(M^n \subset \mathbb{S}^{n+p}(1)\) be a compact submanifold with **parallel mean curvature vector field** (\(\nabla^\perp H = 0\)) and nonzero \(H\), but **nonparallel second fundamental form**. Classical results (Ferus, Naitoh) show that when the second fundamental form is parallel, the submanifold is homogeneous, and its normal holonomy group is the isotropy of an \(s\)-representation. Moreover, Olmos’ normal holonomy theorem implies that if the normal holonomy action is transitive on the unit normal sphere, then \(M\) is umbilical in its normal directions.  

**Open problem:**  
Assume \(M\) is a compact PMC (parallel mean curvature) submanifold in \(\mathbb{S}^{n+p}(1)\) with **irreducible but nontransitive normal holonomy**.  
1. Must \(M\) be **homogeneous**, i.e., an orbit of an isometric group action?  
2. Do such submanifolds actually exist, or does irreducibility of the normal holonomy force parallelism of the full second fundamental form (hence extrinsic symmetry)?  
3. More generally, can one classify compact PMC submanifolds in the sphere whose normal holonomy representation is irreducible?

Why is it a "good" problem:  
This problem probes the exact boundary between **partial extrinsic symmetry** and **full geometric rigidity**.  
- On one side lies the completely classified world of submanifolds with parallel second fundamental form (homogeneous, \(s\)-representation orbits).  
- On the other lies the flexible class of PMC submanifolds, much less understood.  

By injecting the constraint of **irreducible normal holonomy**, the problem cuts to the conceptual core: does the geometric assumption of parallel mean curvature already impose global homogeneity, when coupled with strong holonomy irreducibility?  

It is “good” because:  
- It **builds directly on known rigidity theorems** (Ferus–Naitoh, Olmos) yet asks a question that is not settled by them.  
- It integrates **geometry of mean curvature**, **representation‑theoretic holonomy theory**, and **homogeneous submanifold classification**, bridging multiple areas.  
- It is **simple to state but intricately deep**: tightening the holonomy condition could either uncover new, subtle examples or yield a powerful extension of rigidity results.  
- Either outcome—existence or nonexistence—would advance understanding of how local curvature conditions govern global symmetry.